\documentclass{article}%
\usepackage[T1]{fontenc}%
\usepackage[utf8]{inputenc}%
\usepackage{lmodern}%
\usepackage{textcomp}%
\usepackage{lastpage}%
\usepackage[spanish]{babel}%
\usepackage{array}%
\usepackage{float}% 
\usepackage[letterpaper,top=1.5cm,bottom=1.5cm,left=2.3cm,right=1.5cm]{geometry}%
\usepackage{longtable}%
%
%
%
\begin{document}%
\normalsize%
\section{Lista de Requerimientos}%
\label{sec:ListadeRequerimientos}%
\subsection{Categoría R.SH}%
\label{subsec:CategoraR.SH}%
Requerimientos de Stakeholders (de interesados).\newline%
\newline%
%
\begin{longtable}{|p{2cm}|p{9cm}|p{2.5cm}|}%
\hline%
ID&Descripción del requerimiento&Requisito previo\\%
\hline%
\endhead%
\hline%
\multicolumn{3}{|r|}{Continúa en la siguiente página...}\\%
\hline%
\endfoot%
\hline%
\multicolumn{3}{|r|}{Fin de la tabla}\\%
\hline%
\endlastfoot%
R.SH{-}1& El sistema deberá poder calcular el índice TGBH, acorde a lo establecido en la norma 7243.&{-}\\%
\hline%
R.SH{-}2& El sistema deberá poder calcular el índice ISR (índice de sudoración requerida).&{-}\\%
\hline%
R.SH{-}3& El sistema deberá poder calcular el índice UV.&{-}\\%
\hline%
R.SH{-}4& Los datos medidos e índices de estrés térmico calculados deberán poder ser almacenados en un servidor en la nube.&{-}\\%
\hline%
R.SH{-}5& Los datos deberán poder ser almacenados localmente, al menos por seis meses.&{-}\\%
\hline%
R.SH{-}6& El sistema deberá poder resistir las condiciones ambientales del lugar de instalación, en las fincas de caña de azúcar de Costa Rica, incluyendo polvo y lluvia.&{-}\\%
\hline%
R.SH{-}7& El sistema deberá generar alertas a través del software de monitoreo de condiciones ambientales, cuando los índices de estrés térmico que calcule el sistema, excedan un límite dado.&{-}\\%
\hline%
R.SH{-}8& El sistema deberá poder ser actualizable en su firmware, de manera local.&{-}\\%
\hline%
R.SH{-}9& El sistema deberá cumplir con una verificación con un monitor de TGBH de referencia.&{-}\\%
\hline%
R.SH{-}10& El sistema deberá poder ser transportable en un automóvil o vehículo doble tracción.&{-}\\%
\hline%
R.SH{-}11& El sistema de almacenamiento en nube, cálculos matemáticos, gráficos y alertas deberá poder ser actualizable en su software, de manera remota.&{-}\\%
\hline%
R.SH{-}12& El sistema deberá poder ser alimentado de manera aislada, en caso de no tener acceso a electricidad de la red eléctrica comercial&{-}\\%
\hline%
\end{longtable}

%
\newpage%
\subsection{Categoría R.AVMS}%
\label{subsec:CategoraR.AVMS}%
Requerimientos del Sistema de Medición de Variables Ambientales (AVMS).\newline%
\newline%
%
\begin{longtable}{|p{2cm}|p{9cm}|p{2.7cm}|}%
\hline%
ID&Descripción del requerimiento&Requisito previo\\%
\hline%
\endhead%
\hline%
\multicolumn{3}{|r|}{Continúa en la siguiente página...}\\%
\hline%
\endfoot%
\hline%
\multicolumn{3}{|r|}{Fin de la tabla}\\%
\hline%
\endlastfoot%
R.AVMS{-}1& El sistema deberá poder medir la temperatura seca& R.SH{-}1\\%
\hline%
R.AVMS{-}2& El sistema deberá poder medir la temperatura de globo& R.SH{-}1\\%
\hline%
R.AVMS{-}3& El sistema deberá poder medir la velocidad del viento& R.SH{-}1\\%
\hline%
R.AVMS{-}4& El sistema deberá poder medir la humedad relativa& R.SH{-}1\\%
\hline%
R.AVMS{-}5& El sistema deberá poder medir la radiación UV{-}A& R.SH{-}3\\%
\hline%
R.AVMS{-}6& El sistema deberá poder medir la radiación UV{-}B& R.SH{-}3\\%
\hline%
R.AVMS{-}7& El sistema deberá tener una memoria de almacenamiento masivo local& R.SH{-}5\\%
\hline%
R.AVMS{-}8& El sistema deberá tener una resistencia al polvo&R.SH{-}6\\%
\hline%
R.AVMS{-}9& El sistema deberá permitir su limpieza del polvo durante sus mantenimientos&R.SH{-}6\\%
\hline%
R.AVMS{-}10& El sistema deberá tener una resistencia a la lluvia&R.SH{-}6\\%
\hline%
R.AVMS{-}11& El sistema deberá admitir ser actualizable a través de un protocolo USB&R.SH{-}8\\%
\hline%
R.AVMS{-}12& El sistema deberá tener un error menor al 3 \% al ser comparado con el TGBH de referencia&R.SH{-}9\\%
\hline%
R.AVMS{-}13& El sistema deberá poder ser plegable, lo suficiente para que pueda ser transportable de forma segurda en el techo de un vehículo&R.SH{-}10\\%
\hline%
R.AVMS{-}14& El sistema deberá contar con un sistema de energía renovable&R.SH{-}12\\%
\hline%
R.AVMS{-}15& El sistema deberá contar con un sistema de baterías&R.SH{-}12\\%
\hline%
R.AVMS{-}16& El sistema deberá contar con acceso a internet&R.SH{-}4\\%
\hline%
R.AVMS{-}17& El sistema deberá poder codificar en hexadecimal los datos transmitidos al servidor&R.SH{-}4\\%
\hline%
R.AVMS{-}18& El sistema deberá poder enviar datos a través del protocolo MQTT&R.SH{-}4\\%
\hline%
R.AVMS{-}19& El sistema deberá tener un sistema de protección de sobretensión&R.SH{-}12\\%
\hline%
R.AVMS{-}20& El sistema deberá tener un sistema de estabilización de energía entregada a los componentes electrónicos&R.SH{-}12\\%
\hline%
R.AVMS{-}21& El sistema deberá tener una máquina de estados, que le permita modificar su operación según variables previamente programadas&R.SH{-}12\\%
\hline%
R.AVMS{-}22& El sistema deberá poder soportar temperaturas ambientales de hasta 45 grados Celcius&R.SH{-}6\\%
\hline%
R.AVMS{-}23& El sistema deberá poder soportar humedades relativas desde el 0 \% hasta el 100 \% &R.SH{-}6\\%
\hline%
R.AVMS{-}24& El sistema deberá poder calcular promedios de datos provenientes de variables medidas&R.SH{-}1, R.SH{-}2, R.SH{-}3 y R.SH{-}4\\%
\hline%
\end{longtable}

%
\newpage%
\subsection{Categoría R.PDS}%
\label{subsec:CategoraR.PDS}%
Requerimientos del sistema de potencia y energía (Power Delivery System [PDS]).\newline%
\newline%
%
\begin{longtable}{|p{2cm}|p{9cm}|p{2.5cm}|}%
\hline%
ID&Descripción del requerimiento&Requisito previo\\%
\hline%
\endhead%
\hline%
\multicolumn{3}{|r|}{Continúa en la siguiente página...}\\%
\hline%
\endfoot%
\hline%
\multicolumn{3}{|r|}{Fin de la tabla}\\%
\hline%
\endlastfoot%
R.PDS{-}1& El sistema deberá contar con una fuente de alimentación capaz de entregar voltaje estable de 3.6 V \$\textbackslash{}pm\$+{-} 0.1 V& R.AVMS{-}5\\%
\hline%
R.PDS{-}2& La fuente de alimentación deberá tener protección contra sobrecorriente y cortocircuito& R.AVMS{-}5\\%
\hline%
R.PDS{-}3& El sistema deberá contar con una batería recargable con autonomía mínima de 14 horas& R.AVMS{-}6\\%
\hline%
R.PDS{-}4& La batería deberá tener protección contra descarga profunda y sobrecarga& R.AVMS{-}6\\%
\hline%
R.PDS{-}5& El sistema deberá guardar datos sobre nivel de carga, estado de batería y generación renovable.&R.AVMS{-}14 y R.AVMS{-}15\\%
\hline%
R.PDS{-}6& El sistema deberá tener encapsulado con protección ambiental IP65& R.AVMS{-}10\\%
\hline%
R.PDS{-}7& El sistema deberá permitir recarga mediante fuente solar o alimentador de 12 VDC& R.AVMS{-}6\\%
\hline%
R.PDS{-}8& El sistema de potencia deberá comunicar al procesador central el estado energético para gestión inteligente.& R.AVMS{-}21\\%
\hline%
\end{longtable}

%
\newpage%
\subsection{Categoría R.COMS}%
\label{subsec:CategoraR.COMS}%
Requerimientos de comunicaciones y transmisión de datos [COMS].\newline%
\newline%
%
\begin{longtable}{|p{2cm}|p{9cm}|p{2.5cm}|}%
\hline%
ID&Descripción del requerimiento&Requisito previo\\%
\hline%
\endhead%
\hline%
\multicolumn{3}{|r|}{Continúa en la siguiente página...}\\%
\hline%
\endfoot%
\hline%
\multicolumn{3}{|r|}{Fin de la tabla}\\%
\hline%
\endlastfoot%
R.COMS{-}1& Módulo de red: El sistema deberá contar con un módulo de comunicación WiFi que permita enviar datos al servidor en la nube&R.AVMS{-}16\\%
\hline%
R.COMS{-}2& Módulo de red: El sistema deberá contar con un módulo de comunicación LTE que permita enviar datos al servidor en la nube&R.AVMS{-}16\\%
\hline%
R.COMS{-}3& Codificación de datos: El sistema deberá implementar un algoritmo que convierta los datos en formato hexadecimal antes de la transmisión.&R.AVMS{-}17\\%
\hline%
R.COMS{-}4& Protocolo de transmisión: El sistema deberá soportar MQTT con cifrado TLS para enviar datos de manera segura y confiable.&R.AVMS{-}18\\%
\hline%
R.COMS{-}5& Redundancia de comunicación: El sistema deberá contar con mecanismos de reconexión automática en caso de pérdida de señal.&R.AVMS{-}16\\%
\hline%
R.COMS{-}6& Integridad de datos: El sistema deberá verificar la integridad de los paquetes transmitidos mediante checksum o hash.&R.AVMS{-}17 y R.AVMS{-}18\\%
\hline%
\end{longtable}

%
\newpage%
\subsection{Categoría R.CPU}%
\label{subsec:CategoraR.CPU}%
Requerimientos del procesamiento central y firmware [CPU].\newline%
\newline%
%
\begin{longtable}{|p{2cm}|p{9cm}|p{2.5cm}|}%
\hline%
ID&Descripción del requerimiento&Requisito previo\\%
\hline%
\endhead%
\hline%
\multicolumn{3}{|r|}{Continúa en la siguiente página...}\\%
\hline%
\endfoot%
\hline%
\multicolumn{3}{|r|}{Fin de la tabla}\\%
\hline%
\endlastfoot%
R.CPU{-}1& Procesamiento central: El sistema deberá contar con un procesador capaz de ejecutar cálculos matemáticos en tiempo real&R.AVMS{-}1, R.AVMS{-}2, R.AVMS{-}3 y R.AVMS{-}4\\%
\hline%
R.CPU{-}2& Firmware actualizable: El CPU deberá contar con un conector tipo USB.&R.AVMS{-}11\\%
\hline%
R.CPU{-}3& Validación de sensores: El CPU deberá poder verificar si los sensores lograron conectarse& \\%
\hline%
R.CPU{-}4& Medición sin interrupción: El CPU deberá poder mantener el intervalo de medición, cuando transmita los datos al servidor&R.AVMS{-}4\\%
\hline%
R.CPU{-}5& Gestión de memoria: el CPU deberá administrar almacenamiento temporal (buffer) para evitar pérdida de datos durante transmisión& R.AVMS{-}7\\%
\hline%
R.CPU{-}6& Sincronización temporal: el CPU deberá mantener reloj interno o sincronización con servidor para trazabilidad de mediciones& R.AVMS{-}16\\%
\hline%
R.CPU{-}7& Consumo eficiente: el CPU deberá optimizar su consumo energético para extender la autonomía del sistema, en caso de requerirse& R.AVMS{-}6\\%
\hline%
R.CPU{-}8& Seguridad de firmware: el CPU deberá permitir verificación de integridad del firmware antes de actualización& R.AVMS{-}11\\%
\hline%
\end{longtable}

%
\newpage%
\subsection{Categoría R.DT}%
\label{subsec:CategoraR.DT}%
Requerimientos de sensores de temperatura seca [DT].\newline%
\newline%
%
\begin{longtable}{|p{2cm}|p{9cm}|p{2.5cm}|}%
\hline%
ID&Descripción del requerimiento&Requisito previo\\%
\hline%
\endhead%
\hline%
\multicolumn{3}{|r|}{Continúa en la siguiente página...}\\%
\hline%
\endfoot%
\hline%
\multicolumn{3}{|r|}{Fin de la tabla}\\%
\hline%
\endlastfoot%
R.DT{-}1& La medición de temperatura debe tener una precisión mínima de $\pm$ 0.5 grados Celcius en el rango de 10 grados Celcius a 60 grados Celcius&R.AVMS{-}1\\%
\hline%
R.DT{-}2& La medición de temperatura debe tener una resolución de al menos 0.1 grados Celcius&R.AVMS{-}1\\%
\hline%
R.DT{-}3& El sensor deberá tener un tiempo de respuesta térmica menor a 10 segundos& R.AVMS{-}1\\%
\hline%
R.DT{-}4& El sensor deberá estar calibrado contra un patrón trazable a normas internacionales (ISO/NIST)& R.AVMS{-}1\\%
\hline%
R.DT{-}5& El sensor deberá contar con protección contra radiación solar directa (pantalla o encapsulado)& R.AVMS{-}1\\%
\hline%
R.DT{-}6& El sensor deberá ser compatible con el sistema de adquisición de datos del CPU& R.AVMS{-}1, R.CPU{-}3\\%
\hline%
R.DT{-}7& El sensor deberá tener protección ambiental contra lluvia y polvo& R.AVMS{-}8, R.AVMS{-}10\\%
\hline%
\end{longtable}

%
\newpage%
\subsection{Categoría R.GB}%
\label{subsec:CategoraR.GB}%
Requerimientos de sensores de temperatura de globo [GB].\newline%
\newline%
%
\begin{longtable}{|p{2cm}|p{9cm}|p{2.5cm}|}%
\hline%
ID&Descripción del requerimiento&Requisito previo\\%
\hline%
\endhead%
\hline%
\multicolumn{3}{|r|}{Continúa en la siguiente página...}\\%
\hline%
\endfoot%
\hline%
\multicolumn{3}{|r|}{Fin de la tabla}\\%
\hline%
\endlastfoot%
R.GB{-}1& La medición de temperatura de globo deberá tener una precisión mínima de $\pm$ 0.5 grados Celcius en el rango de 10 \$\^{}\textbackslash{}\{\textbackslash{}circ\textbackslash{}\}\$C a 60 \$\^{}\textbackslash{}\{\textbackslash{}circ\textbackslash{}\}\$C& R.AVMS{-}2\\%
\hline%
R.GB{-}2& La medición de temperatura de globo deberá tener una resolución de al menos 0.1 grados Celcius& R.AVMS{-}2\\%
\hline%
R.GB{-}3& El sensor deberá tener un tiempo de respuesta térmica menor a 20 segundos& R.AVMS{-}2\\%
\hline%
R.GB{-}4& El sensor deberá estar calibrado contra un patrón trazable a normas internacionales (ISO/NIST)& R.AVMS{-}2\\%
\hline%
R.GB{-}5& El sensor deberá estar montado dentro de un globo negro, acorde a ISO 7243& R.AVMS{-}2\\%
\hline%
R.GB{-}6& El sensor deberá ser compatible con el sistema de adquisición de datos del CPU& R.AVMS{-}2, R.CPU{-}3\\%
\hline%
R.GB{-}7& El sensor deberá tener protección ambiental contra polvo y lluvia& R.AVMS{-}8, R.AVMS{-}10\\%
\hline%
\end{longtable}

%
\newpage%
\subsection{Categoría R.UV}%
\label{subsec:CategoraR.UV}%
Requerimientos de sensores de radiación ultravioleta [UV].\newline%
\newline%
%
\begin{longtable}{|p{2cm}|p{9cm}|p{2.5cm}|}%
\hline%
ID&Descripción del requerimiento&Requisito previo\\%
\hline%
\endhead%
\hline%
\multicolumn{3}{|r|}{Continúa en la siguiente página...}\\%
\hline%
\endfoot%
\hline%
\multicolumn{3}{|r|}{Fin de la tabla}\\%
\hline%
\endlastfoot%
R.UV{-}1& El sensor deberá medir radiación ultravioleta en el rango de 280 nm a 400 nm (UV{-}A y UV{-}B)& R.AVMS{-}3\\%
\hline%
R.UV{-}2& La medición de radiación UV deberá tener una precisión mínima de $\pm$ 5 \% respecto a un patrón de referencia& R.AVMS{-}3\\%
\hline%
R.UV{-}3& El sensor deberá tener una resolución mínima de 0.1 mW/cm$^2$ & R.AVMS{-}3\\%
\hline%
R.UV{-}4& El sensor deberá estar calibrado contra un patrón trazable a normas internacionales (ISO/NIST)& R.AVMS{-}3\\%
\hline%
R.UV{-}5& El sensor deberá contar con protección óptica para evitar saturación por radiación solar directa& R.AVMS{-}3\\%
\hline%
R.UV{-}6& El sensor deberá ser compatible con el sistema de adquisición de datos del CPU& R.AVMS{-}3, R.CPU{-}3\\%
\hline%
R.UV{-}7& El sensor deberá tener encapsulado con protección ambiental contra polvo y lluvia& R.AVMS{-}8, R.AVMS{-}10\\%
\hline%
\end{longtable}

%
\newpage%
\subsection{Categoría R.WS}%
\label{subsec:CategoraR.WS}%
Requerimientos de sensor de velocidad del viento [WS].\newline%
\newline%
%
\begin{longtable}{|p{2cm}|p{9cm}|p{2.5cm}|}%
\hline%
ID&Descripción del requerimiento&Requisito previo\\%
\hline%
\endhead%
\hline%
\multicolumn{3}{|r|}{Continúa en la siguiente página...}\\%
\hline%
\endfoot%
\hline%
\multicolumn{3}{|r|}{Fin de la tabla}\\%
\hline%
\endlastfoot%
R.WS{-}1& El sensor deberá medir velocidad del viento en el rango de 0.3 m/s a 20 m/s& R.AVMS{-}3\\%
\hline%
R.WS{-}2& La medición de velocidad del viento deberá tener una precisión mínima de +{-} 0.3 m/s& R.AVMS{-}3\\%
\hline%
R.WS{-}3& El sensor deberá tener una resolución mínima de 0.1 m/s& R.AVMS{-}3\\%
\hline%
R.WS{-}4& El sensor deberá estar calibrado contra un patrón trazable a normas internacionales (ISO/NIST)& R.AVMS{-}3\\%
\hline%
R.WS{-}5& El sensor deberá tener un tiempo de respuesta menor a 30 segundos& R.AVMS{-}3\\%
\hline%
R.WS{-}6& El sensor deberá ser compatible con el sistema de adquisición de datos del CPU& R.AVMS{-}3, R.CPU{-}3\\%
\hline%
R.WS{-}7& El sensor deberá tener encapsulado con protección ambiental contra polvo e insectos& R.AVMS{-}8, R.AVMS{-}10\\%
\hline%
R.WS{-}8& El sensor deberá contar con carcasa protectora y drenaje que asegure resistencia a lluvia y condensación, manteniendo la precisión de la medición& R.AVMS{-}10\\%
\hline%
\end{longtable}

%
\newpage%
\subsection{Categoría R.STR}%
\label{subsec:CategoraR.STR}%
Requerimientos de la estructura física del sistema [STR].\newline%
\newline%
%
\begin{longtable}{|p{2cm}|p{9cm}|p{2.5cm}|}%
\hline%
ID&Descripción del requerimiento&Requisito previo\\%
\hline%
\endhead%
\hline%
\multicolumn{3}{|r|}{Continúa en la siguiente página...}\\%
\hline%
\endfoot%
\hline%
\multicolumn{3}{|r|}{Fin de la tabla}\\%
\hline%
\endlastfoot%
R.STR{-}1& La estructura deberá permitir instalar los sensores de medición a la altura especificada por ISO 7243 (1.1 m a 1.7 m sobre el suelo)& R.AVMS{-}12\\%
\hline%
R.STR{-}2& La estructura deberá contar con soporte superior para el sensor UV y el panel solar& R.AVMS{-}12, R.AVMS{-}6\\%
\hline%
R.STR{-}3& La estructura deberá ser resistente a condiciones ambientales exteriores (humedad, lluvia, viento)& R.AVMS{-}10\\%
\hline%
R.STR{-}4& La estructura deberá estar fabricada con materiales anticorrosivos (ej. aluminio anodizado, acero inoxidable, polímeros técnicos)& R.AVMS{-}10\\%
\hline%
R.STR{-}5& La estructura deberá tener estabilidad mecánica suficiente para soportar ráfagas de viento de hasta 20 m/s sin comprometer la medición& R.AVMS{-}3, R.AVMS{-}10\\%
\hline%
R.STR{-}6& La estructura deberá permitir montaje y desmontaje modular para facilitar transporte y mantenimiento& R.AVMS{-}12\\%
\hline%
R.STR{-}7& La estructura deberá contar con protección IP65 en puntos de conexión eléctrica (panel solar, CPU, sensores)& R.AVMS{-}8, R.AVMS{-}10\\%
\hline%
\end{longtable}

%
\newpage%
\subsection{Categoría R.SDPS}%
\label{subsec:CategoraR.SDPS}%
Requerimientos del servidor en la nube y procesamiento de datos [SDPS].\newline%
\newline%
%
\begin{longtable}{|p{2cm}|p{9cm}|p{2.5cm}|}%
\hline%
ID&Descripción del requerimiento&Requisito previo\\%
\hline%
\endhead%
\hline%
\multicolumn{3}{|r|}{Continúa en la siguiente página...}\\%
\hline%
\endfoot%
\hline%
\multicolumn{3}{|r|}{Fin de la tabla}\\%
\hline%
\endlastfoot%
R.SDPS{-}1&El servidor en la nube deberá poder manejar múltiples dispositivos y grandes volúmenes de datos sin degradar el rendimiento.&R.SH{-}4\\%
\hline%
R.SDPS{-}2&Los datos almacenados deberán ser accesibles desde cualquier ubicación autorizada, con credenciales de usuario&R.SH{-}4\\%
\hline%
R.SDPS{-}3&La base de datos en la nube deberá poder recibir datos de sistemas embebidos IoT a través de MQTT&R.SH{-}4\\%
\hline%
R.SDPS{-}4&La base de datos en la nube deberá poder enviar datos cuando un servidor lo solicite&R.SH{-}4\\%
\hline%
R.SDPS{-}6&El sistema deberá poder realizar cálculos matemáticos con ecuaciones polinomiales de varias variables& R.SH{-}1, R.SH{-}2, R.SH{-}3\\%
\hline%
R.SDPS{-}7&El sistema deberá poder almacenar los resultados de las operaciones matemáticas realizadas, asociadas con los resultados de los índices de estrés térmico& R.SH{-}1, R.SH{-}2, R.SH{-}3\\%
\hline%
\end{longtable}

%
\newpage%
\subsection{Categoría R.UICS}%
\label{subsec:CategoraR.UICS}%
Requerimientos de la interfaz de usuario y alertas [UICS].\newline%
\newline%
%
\begin{longtable}{|p{2cm}|p{9cm}|p{2.5cm}|}%
\hline%
ID&Descripción del requerimiento&Requisito previo\\%
\hline%
\endhead%
\hline%
\multicolumn{3}{|r|}{Continúa en la siguiente página...}\\%
\hline%
\endfoot%
\hline%
\multicolumn{3}{|r|}{Fin de la tabla}\\%
\hline%
\endlastfoot%
R.UICS{-}1& El sistema deberá generar alertas a través del software de monitoreo de condiciones ambientales, cuando los índices de estrés térmico que calcule el sistema, excedan un límite dado.&R.SH{-}7\\%
\hline%
R.UICS{-}2& El sistema deberá comparar continuamente los valores calculados con los umbrales definidos.&R.SH{-}7\\%
\hline%
R.UICS{-}3& El sistema deberá emitir una alerta cuando se exceda un límite, en tiempo real.&R.SH{-}7\\%
\hline%
R.UICS{-}4& El sistema deberá soportar diferentes modalidades de alerta (visual en pantalla, sonora, notificación en la nube y por correo).&R.SH{-}7\\%
\hline%
R.UICS{-}5& El sistema deberá almacenar un historial de alertas generadas, con fecha, hora y valor del índice.&R.SH{-}7\\%
\hline%
R.UICS{-}6& El sistema deberá permitir ajustar la severidad de las alertas (ej. advertencia, crítico).&R.SH{-}11\\%
\hline%
R.UICS{-}7& El sistema deberá permitir la actualización del software desde ubicaciones remotas, con protocolos de seguridad (autenticación, cifrado).&R.SH{-}11\\%
\hline%
R.UICS{-}8& El sistema deberá garantizar que las actualizaciones no afecten la integridad de los datos almacenados ni la configuración previa.&R.SH{-}11\\%
\hline%
R.UICS{-}9& El sistema deberá contar con un procedimiento para revertir la actualización en caso de fallo.&R.SH{-}11\\%
\hline%
R.UICS{-}10& El sistema deberá informar al usuario cuando una actualización esté disponible y cuando se haya completado.&R.SH{-}11\\%
\hline%
R.UICS{-}11& El sistema deberá minimizar el tiempo de inactividad durante el proceso de actualización.&R.SH{-}11\\%
\hline%
R.UICS{-}12& El sistema deberá mantener un historial de versiones del software y actualizaciones aplicadas.&\\%
\hline%
\end{longtable}

%
\newpage%
\subsection{Categoría  }%
\label{subsec:Categora}%
Descripción no definida para esta categoría.\newline%
\newline%
%
\begin{longtable}{|p{2cm}|p{9cm}|p{2.5cm}|}%
\hline%
ID&Descripción del requerimiento&Requisito previo\\%
\hline%
\endhead%
\hline%
\multicolumn{3}{|r|}{Continúa en la siguiente página...}\\%
\hline%
\endfoot%
\hline%
\multicolumn{3}{|r|}{Fin de la tabla}\\%
\hline%
\endlastfoot%
 &&\\%
\hline%
\end{longtable}


%
\end{document}